% !TeX root = ../main.tex
% chapters/results.tex

\chapter{Results}
\label{ch:results}

This chapter presents the main results of the formalisation.
\Cref{sec:res-overview} summarises the library's scope and
metrics.  \Cref{sec:res-soundness,sec:res-adequacy} state and
prove the two central metatheorems---soundness and adequacy---for the
PCF$_v$ language, and \cref{sec:res-correspondence} combines them into a
single convergence correspondence.
\Cref{sec:res-enriched} records the enriched-category results:
the Yoneda lemma and recursive domain equations.
Finally, \cref{sec:res-examples} presents worked examples that
exercise every layer of the library.

% ── Library metrics ───────────────────────────────────────────────
\section{Library Overview}
\label{sec:res-overview}

\Cref{tab:library-metrics} summarises the size of the project.
The library currently comprises 35 source files across five directories,
totalling approximately 15{,}700 lines of Rocq code; the remaining
lines belong to the four example walkthroughs and the seven
test suites.

\begin{table}[ht]
\centering
\begin{tabular}{l r r l}
\hline
\textbf{Layer} & \textbf{Files} & \textbf{Lines} & \textbf{Contents} \\
\hline
\texttt{structures/} &  4 &    974 & \HB{} mixins and structures \\
\texttt{theory/}     & 13 &  7{,}488 & Core domain theory \\
\texttt{instances/}  &  6 &  2{,}084 & Concrete CPO instances \\
\texttt{lang/}       &  7 &  3{,}402 & PCF \& QMiniCore \\
\texttt{quantum/}    &  5 &  1{,}767 & Quantum CPO layer \\
\hline
\emph{Library sub-total} & 35 & 15{,}715 & \\
\hline
\texttt{examples/}   &  4 &    986 & Worked walkthroughs \\
\texttt{test/}       &  7 &  1{,}579 & Property-level test suites \\
\hline
\textbf{Total}       & \textbf{46} & \textbf{18{,}280} & \\
\hline
\end{tabular}
\caption{Library metrics by directory layer.  Line counts are over
  \texttt{.v} source files, not including generated build
  artifacts.}
\label{tab:library-metrics}
\end{table}

Across the library, the code base contains approximately 240
definitions and 601 lemmas or theorems.
The \HB{} hierarchy declares 25 structures and 81 instances,
spanning the full stack from \texttt{PartialOrder} up to
\texttt{MixedLCFunctor}.

The core library (the first five rows of \cref{tab:library-metrics})
contains \textbf{zero} uses of \texttt{Admitted}.  One explicitly
declared axiom exists: \texttt{bilimit\_exists} in
\texttt{DomainEquations.v}, which postulates the existence of a
bilimit for any approximation sequence
(\cref{sec:res-domain-equations}).  The experimental
\texttt{LiftMonad.v} file (the coinductive lift monad explored in
\cref{sec:disc-lift-monad}) contains two \texttt{Admitted} lemmas,
both blocked by the guardedness obstruction discussed there; this file
is not imported by any other module.

The library builds under \Rocq~9.0 with the Hierarchy Builder
plugin, using the Dune build system.  All 46 files compile
without warnings.

% ── Soundness ─────────────────────────────────────────────────────
\section{PCF Soundness}
\label{sec:res-soundness}

% Source: lang/PCF_Soundness.v

The mathematical proof of soundness was given in \cref{sec:bg-adequacy}.
Here we present the mechanised version in \Rocq{} and note how the
formal proof mirrors the paper argument.

\begin{theorem}[Mechanised soundness]
\label{thm:mech-soundness}
For every closed expression~$e$ and closed value~$v$,
\[
  e \Downarrow v
  \;\implies\;
  \texttt{sem\_exp}\;e\;\mathtt{tt}
    = \mathsf{Some}\,(\texttt{sem\_val}\;v\;\mathtt{tt}).
\]
\end{theorem}

\noindent
The complete proof script is shown in \cref{lst:soundness}.

\begin{lstlisting}[caption={Soundness theorem and proof (\texttt{PCF\_Soundness.v}).},
                    label={lst:soundness}]
Theorem soundness :
    forall {τ : Ty} (e : CExp τ) (v : CValue τ),
      e ⇓ v  →  sem_exp e tt = Some (sem_val v tt).
Proof.
  intros τ e v H.
  induction H as
    [ t' v'                                     (* e_Val   *)
    | t1 t2 e1 w1 e2 w2  He1 IH1  He2 IH2      (* e_Let   *)
    | t1 t2 body varg w  He IH                   (* e_App   *)
    | t1 t2 v1 v2                                (* e_Fst   *)
    | t1 t2 v1 v2                                (* e_Snd   *)
    | op n1 n2                                    (* e_Op    *)
    | n1 n2                                       (* e_Gt    *)
    | t' e1 e2 w  He IH                          (* e_IfTrue  *)
    | t' e1 e2 w  He IH                          (* e_IfFalse *)
    ].
  - (* e_Val *)  reflexivity.
  - (* e_Let *)  rewrite sem_exp_LET, IH1.
                 rewrite <- IH2, sem_single_subst.
                 reflexivity.
  - (* e_App *)  rewrite sem_exp_APP, sem_val_FIX_unfold.
                 rewrite <- IH, sem_double_subst.
                 reflexivity.
  - (* e_Fst *)  rewrite sem_exp_FST, sem_val_PAIR. reflexivity.
  - (* e_Snd *)  rewrite sem_exp_SND, sem_val_PAIR. reflexivity.
  - (* e_Op  *)  rewrite sem_exp_OP. reflexivity.
  - (* e_Gt  *)  rewrite sem_exp_GT. reflexivity.
  - (* e_IfTrue  *) rewrite sem_exp_IFB. simpl. exact IH.
  - (* e_IfFalse *) rewrite sem_exp_IFB. simpl. exact IH.
Qed.
\end{lstlisting}

The proof proceeds by induction on the evaluation derivation
\texttt{H~:~e~$\Downarrow$~v}, exactly as in the paper proof
(\cref{thm:soundness}).  Each case applies the $\beta$-computation
rule exported by \texttt{PCF\_Denotational.v}---for example,
\texttt{sem\_exp\_LET} for the \textsc{E-Let} rule and
\texttt{sem\_exp\_APP} together with \texttt{sem\_val\_FIX\_unfold}
for the \textsc{E-App} rule---followed by the induction hypothesis.
The two non-trivial binder cases (\textsc{E-Let} and \textsc{E-App})
additionally appeal to the substitution corollaries
\texttt{sem\_single\_subst} and \texttt{sem\_double\_subst}
developed in \cref{sec:meth-pcf}.

The short proof scripts reflect a general pattern: because the
denotational semantics was defined to compute definitionally on each
syntactic form (DD-014), each case reduces to a chain of
\texttt{rewrite} steps that reassemble the denotational equations
\emph{without case analysis on the domain elements themselves}.
The intrinsically typed representation
(\cref{sec:meth-pcf}) eliminates all variable-binding boilerplate:
there are no well-scopedness side-conditions to discharge.

\medskip

\noindent
Three corollaries follow directly:

\begin{corollary}[Convergence implies definedness]
\label{cor:sound-not-none}
If\/ $e \Downarrow$ then $\sem{e}^c_E \neq \mathsf{None}$.
\end{corollary}

\begin{corollary}[Values denote themselves]
\label{cor:sound-val}
$\texttt{sem\_exp}\;(\mathsf{val}\;v)\;\mathtt{tt}
  = \mathsf{Some}\,(\texttt{sem\_val}\;v\;\mathtt{tt})$.
\end{corollary}

\begin{corollary}[Denotational agreement]
\label{cor:sound-agree}
If\/ $e \Downarrow v_1$ and $e \Downarrow v_2$ then
$\texttt{sem\_val}\;v_1\;\mathtt{tt}
  = \texttt{sem\_val}\;v_2\;\mathtt{tt}$.
\end{corollary}

\noindent
\Cref{cor:sound-not-none} provides the ``easy'' direction of the
convergence correspondence; the harder converse is the subject of
the next section.

% ── Adequacy ──────────────────────────────────────────────────────
\section{PCF Adequacy}
\label{sec:res-adequacy}

% Source: lang/PCF_Adequacy.v

The mathematical treatment of adequacy---the logical relation, its
admissibility, the fundamental lemma, and the adequacy
corollary---was given in \cref{sec:bg-adequacy}.  This section presents
the mechanised development in \texttt{PCF\_Adequacy.v} (820~lines) and
highlights the points where the formal proof diverges from, or
elaborates on, the paper argument.

\subsection*{The logical relation}

The value relation \texttt{rel\_val} and expression relation
\texttt{rel\_exp} are defined exactly as in
\cref{def:logical-relation}.  In \Rocq{}, the value relation is a
\texttt{Fixpoint} on the type index~$\tau$, so that the recursive
occurrences in the arrow and product cases are structurally guarded:

\begin{lstlisting}[caption={The logical relation (\texttt{PCF\_Adequacy.v}).},
                    label={lst:rel-val}]
Fixpoint rel_val (τ : Ty) : sem_ty τ → CValue τ → Prop :=
  match τ return sem_ty τ → CValue τ → Prop with
  | Nat    => fun d v => v = NLIT d
  | Bool   => fun b v => v = BLIT b
  | Arrow t1 t2 =>
      fun f v =>
        exists body, v = FIX t1 t2 body /\
        forall d1 v1, rel_val t1 d1 v1 →
          match (f d1) with
          | None   => True
          | Some d => exists w, doubleSubstExp v1 v body ⇓ w
                                /\ rel_val t2 d w
          end
  | Prod t1 t2 =>
      fun '(d1, d2) v =>
        exists v1 v2, v = PAIR v1 v2 /\
          rel_val t1 d1 v1 /\ rel_val t2 d2 v2
  end.

Definition rel_exp (τ : Ty) (o : option (sem_ty τ))
    (e : CExp τ) : Prop :=
  match o with
  | None   => True
  | Some d => exists v, e ⇓ v /\ rel_val τ d v
  end.
\end{lstlisting}

\noindent
The expression relation matches the ``lifted'' formulation of BKV: if
the denotation is $\bot$ (\texttt{None}), the relation is trivially
true; if the denotation is defined (\texttt{Some~d}), the expression
must converge to a value related by \texttt{rel\_val}.  This
formulation avoids the implication-based definition of Plotkin
\cite{Plotkin1977} and simplifies the chain reasoning in the
admissibility proof.

\subsection*{Admissibility}

The admissibility lemmas establish that the logical relation is
compatible with suprema of chains, which is the precondition for
applying fixed-point induction in the \textsc{Fix} case
(\cref{thm:scott-induction}).

\begin{theorem}[Mechanised admissibility]
\label{thm:mech-admissible}
For every type $\tau$ and closed value~$v$, the predicate
$d \mapsto \texttt{rel\_val}\;\tau\;d\;v$ is admissible.
\end{theorem}

\noindent
The proof proceeds by induction on $\tau$, mirroring
\cref{lem:rel-val-admissible}.  In the \Rocq{} formalisation, the
arrow case is the most technically involved: given a chain of
functions $(f_n)$ in the lifted function space
$[\sem{\tau_1} \to_c \sem{\tau_2}_\bot]$, the proof must:
\begin{enumerate}
  \item apply the library's \texttt{chain\_some\_stable} lemma to
    show that once $f_N(d_1) \neq \mathsf{None}$, all later chain
    elements remain defined;
  \item invoke \texttt{eval\_deterministic} to show that the
    convergence witness~$w$ is the same at every index; and
  \item extract a chain in the inner type~$\sem{\tau_2}$ via
    the \texttt{D\_chain} construction from \texttt{Lift.v}, then
    apply the induction hypothesis at $\tau_2$.
\end{enumerate}
Two derived lemmas complete the admissibility interface:
\texttt{rel\_exp\_admissible} lifts the result from values to
expressions, and \texttt{rel\_exp\_admissible\_pointwise} packages it
in the exact form required by the \textsc{Fix} case---as an
admissible predicate on functions
$h \mapsto \forall\, (d_1, v_1),\;
\texttt{rel\_val}\;\tau_1\;d_1\;v_1 \to
\texttt{rel\_exp}\;\tau_2\;(h\;d_1)\;\ldots$

\subsection*{The fundamental lemma}

The fundamental lemma states that every well-typed term, instantiated
by a related environment, lies in the logical relation.  Because
\texttt{Value} and \texttt{Exp} are mutually inductive, the proof
uses a combined induction scheme:

\begin{lstlisting}[caption={Fundamental lemma statement
  (\texttt{PCF\_Adequacy.v}).},
                    label={lst:fundamental}]
Scheme value_ind2 := Induction for Value Sort Prop
  with exp_ind2   := Induction for Exp   Sort Prop.
Combined Scheme val_exp_ind from value_ind2, exp_ind2.

Theorem fundamental_lemma :
    (forall {Γ τ} (v : Value Γ τ) σ γ,
       rel_env Γ γ σ →
       rel_val τ (sem_val v γ) (substVal σ v))
  /\
    (forall {Γ τ} (e : Exp Γ τ) σ γ,
       rel_env Γ γ σ →
       rel_exp τ (sem_exp e γ) (substExp σ e)).
\end{lstlisting}

\noindent
Most cases mirror the paper proof in \cref{thm:fundamental-lemma}
directly: variables appeal to the environment relation, literals are
immediate, pairs use the induction hypothesis on each component, and
expression forms (\textsc{Let}, \textsc{App}, \textsc{IfB}, etc.)
unfold the denotation, case-split on partiality, and assemble the
evaluation derivation from the sub-results.

\medskip

\noindent
\textbf{The \textsc{Fix} case.}\;
This is the most involved case and warrants detailed discussion.
Given $v = \mathsf{fix}\;\tau_1\;\tau_2.\,e$ with denotation
$\sem{v}_V(\gamma) = \fixp(F)$ where
$F(h) = \lambda a.\, \sem{e}_E(a, h, \gamma)$, the goal is to show
that $\fixp(F)$ is related to the substituted value.  Rather than
invoking \texttt{fixp\_ind} directly (which would require matching the
\HB{} \texttt{PointedCPO} instance on the function space), the proof
reasons about the Kleene iterates $F^n(\bot)$ directly:
\begin{enumerate}
  \item A secondary induction on $n : \mathbb{N}$ establishes
    $\texttt{rel\_exp}\;\tau_2\;(F^n(\bot)(d_1))\;
    (\texttt{doubleSubstExp}\;v_1\;v'\;\mathit{body'})$
    for every related argument $(d_1, v_1)$.
    The base case $n = 0$ is trivial ($F^0(\bot) = \bot$, so
    $\texttt{rel\_exp}$ holds vacuously).  The step case constructs
    an extended environment
    $(d_1,\, F^n(\bot),\, \gamma)$ related to the extended substitution
    $[v_1/0,\, v'/1,\, \sigma]$ (using the inner induction hypothesis
    to show $F^n(\bot)$ is related to $v'$), then applies the
    outer induction hypothesis on the body~$e$.
  \item The substitution-composition lemma
    \texttt{subst\_double\_ext} bridges the syntactic substitution
    $[v_1/0,\, v'/1] \circ \sigma^{\uparrow\uparrow}$
    to the operational substitution
    $\texttt{doubleSubstExp}\;v_1\;v'\;\mathit{body'}$.
  \item \texttt{rel\_exp\_admissible\_pointwise} closes the gap from the
    iterates to the fixed point: since the predicate holds at every
    $F^n(\bot)$ and is admissible, it holds at
    $\fixp(F) = \bigsqcup_n F^n(\bot)$.
\end{enumerate}

\subsection*{Adequacy}

Given the fundamental lemma, the adequacy theorem is a short
corollary.

\begin{theorem}[Mechanised adequacy]
\label{thm:mech-adequacy}
For every closed expression~$e$,
\[
  \texttt{sem\_exp}\;e\;\mathtt{tt} \neq \mathsf{None}
  \;\implies\;  e \Downarrow.
\]
\end{theorem}

\noindent
The proof instantiates the fundamental lemma at the empty environment
($\Gamma = \cdot$, $\gamma = \mathtt{tt}$, $\sigma = \mathsf{id}$),
rewrites the trivial identity substitution away with
\texttt{substExp\_id}, and extracts the convergence witness from the
expression relation.  The complete proof occupies eight lines:

\begin{lstlisting}[caption={Adequacy theorem (\texttt{PCF\_Adequacy.v}).},
                    label={lst:adequacy}]
Theorem adequacy :
    forall {τ : Ty} (e : CExp τ),
      sem_exp e tt <> None → e ⇓.
Proof.
  intros τ e Hne.
  destruct fundamental_lemma as [_ Hexp].
  assert (Henv : rel_env [] tt subst_id)
    by (intros τ' x; dependent destruction x).
  specialize (Hexp [] τ e subst_id tt Henv).
  rewrite substExp_id in Hexp.
  destruct (sem_exp e tt) as [d|]; [|contradiction].
  exact (rel_exp_some_converges d e Hexp).
Qed.
\end{lstlisting}

\noindent
The contrast with the soundness proof (\cref{lst:soundness}) is
instructive: soundness is a straightforward induction on evaluation
derivations, while adequacy requires the entire logical-relation
machinery---820~lines of infrastructure culminating in the eight-line
extraction above.

\section{Convergence Correspondence}
\label{sec:res-correspondence}

Combining the soundness theorem (\cref{thm:mech-soundness}) with the
adequacy theorem (\cref{thm:mech-adequacy}) yields the central
metatheoretic result of the PCF formalisation.

\begin{theorem}[Convergence correspondence]
\label{thm:mech-convergence}
For every closed expression $e : \mathsf{CExp}\;\tau$,
\[
  e \Downarrow
  \quad\iff\quad
  \texttt{sem\_exp}\;e\;\mathtt{tt} \neq \mathsf{None}.
\]
\end{theorem}

\noindent
The \Rocq{} proof is a direct assembly of the two directions:

\begin{lstlisting}[caption={Convergence correspondence
  (\texttt{PCF\_Adequacy.v}).},
                    label={lst:convergence}]
Corollary convergence_iff_defined :
    forall {τ : Ty} (e : CExp τ),
      e ⇓  <->  sem_exp e tt <> None.
Proof.
  intros τ e. split.
  - intros [v Hv].
    rewrite (soundness e v Hv). discriminate.
  - exact (adequacy e).
Qed.
\end{lstlisting}

\noindent
The forward direction unfolds the convergence witness to obtain a value
$v$ with $e \Downarrow v$, applies soundness to rewrite the denotation
to $\mathsf{Some}\,(\texttt{sem\_val}\;v\;\mathtt{tt})$, then
discharges the $\neq \mathsf{None}$ goal by discrimination.  The
converse is the adequacy theorem itself.

\medskip

\noindent
\textbf{Interpretation.}\;
The convergence correspondence establishes that the denotational
semantics is a \emph{faithful} account of operational convergence
behaviour: a closed program converges operationally if and only if its
denotation is non-$\bot$.  In particular, no closed program is
``accidentally convergent'' (converging operationally but denoting
$\bot$) or ``accidentally divergent'' (denoting a proper value but
failing to evaluate).

\medskip

\noindent
\textbf{Relation to full abstraction.}\;
The convergence correspondence is strictly weaker than \emph{full
abstraction}, which would require
\[
  \sem{e_1}^c_E = \sem{e_2}^c_E
  \quad\iff\quad
  e_1 \cong_{\mathit{obs}} e_2
\]
for all closed terms of the same type, where $\cong_{\mathit{obs}}$
denotes observational equivalence.  Full abstraction fails for the
standard CPO model of PCF: the denotational semantics distinguishes
too few programs, because the model contains elements---such as the
parallel-or function---that are not the denotation of any PCF
term~\cite{Plotkin1977}.  Our result makes no claim about
observational equivalence; it characterises only the
convergence/divergence distinction.

% ── Enriched category results ─────────────────────────────────────
\section{Enriched Category Results}
\label{sec:res-enriched}

The PCF metatheory of the preceding sections exercises the library's
CPO constructions---products, lifting, function spaces, and fixed
points.  This section turns to the enriched categorical layer:
the Yoneda lemma (\cref{sec:res-yoneda}) and recursive domain
equations (\cref{sec:res-domain-equations}).

\subsection{Yoneda Lemma}
\label{sec:res-yoneda}

% Source: instances/Yoneda.v (~460 lines)

The Yoneda lemma is one of the most fundamental results in category
theory.  In its classical form it characterises the position of an
object within a category entirely through the morphisms into (or out of)
that object: every ``natural'' way of transforming those morphisms is
determined by a single element.

More precisely, fix a locally small category~$\mathcal{C}$, an
object~$X$, and a functor $F : \mathcal{C} \to \mathbf{Set}$.  The
\emph{representable functor} $\mathrm{Hom}(X, -)$ sends each object~$A$
to the set $\mathcal{C}(X, A)$ of morphisms from~$X$ to~$A$, and each
morphism $g : A \to B$ to post-composition $f \mapsto g \circ f$.  The
Yoneda lemma then states:

\begin{theorem}[Yoneda lemma --- set-level]
\label{thm:yoneda-set}
There is a bijection
\[
  \mathrm{Nat}(\mathrm{Hom}(X,-),\; F)
    \;\cong\;
  F(X),
\]
natural in both $X$ and $F$.  The bijection sends a natural
transformation $\alpha : \mathrm{Hom}(X,-) \Rightarrow F$ to the element
$\alpha_X(\mathrm{id}_X) \in F(X)$, and an element $x \in F(X)$ to the
natural transformation with components $\alpha_A(f) = F(f)(x)$.
\end{theorem}

\noindent
When $\mathcal{C}$ is enriched over a monoidal category~$\mathcal{V}$,
the hom-sets become hom-objects in~$\mathcal{V}$ and the bijection
should be promoted to an isomorphism in~$\mathcal{V}$
(Kelly~\cite{Kelly1982}, Ch.~2.4).
In our setting $\mathcal{V} = \CPO$ and
$\mathcal{C} = \CPO$ (self-enriched via \texttt{Function.v}), so the
hom-objects are the continuous function-space CPOs~$[X \to_c A]$ of
\cref{prop:function-space}, and the Yoneda bijection becomes an
isomorphism of CPOs.  The library's \texttt{Yoneda.v} mechanises this
enriched form for fixed~$X$ and~$F$; the parametric naturality in $X$
and $F$ is not yet formalised.

\paragraph{The representable functor.}
For a fixed CPO~$X$, the representable functor
$\mathrm{Hom}(X, -)$ is packaged as an \texttt{lc\_functor}
(\cref{sec:meth-lc-functors}) named \texttt{repr\_functor}: it sends
an object~$D$ to $[X \to_c D]$ and a morphism~$f$ to
post-composition $g \mapsto f \circ g$.  The proofs that
post-composition preserves identities, composition, and $\omega$-chain
suprema are all discharged by \texttt{reflexivity}---each reduces to a
definitional equality.

\paragraph{The enriched Yoneda isomorphism.}
The Yoneda isomorphism of \cref{thm:yoneda-set}, promoted to the
$\CPO$-enriched setting, is witnessed by two maps:
\begin{itemize}
  \item \textbf{Evaluation.}\;
    $\texttt{yoneda\_eval}(\alpha) = \alpha_X(\mathrm{id}_X)$
    extracts the element determined by the natural transformation.
  \item \textbf{Embedding.}\;
    $\texttt{yoneda\_embed}(x)$ constructs the natural transformation
    whose $A$-component sends $f \in [X \to_c A]$ to $F(f)(x)$.
    Naturality follows from the functoriality of~$F$: for
    $g : A \to B$, the naturality square requires
    $F(g)(F(f)(x)) = F(g \circ f)(x)$, which is the
    composition law of~$F$.
\end{itemize}
These maps are mutually inverse.  The left inverse
$\texttt{yoneda\_eval}(\texttt{yoneda\_embed}(x)) = x$ reduces to
$F(\mathrm{id}_X)(x) = x$ by the functor identity law.  The right
inverse $\texttt{yoneda\_embed}(\texttt{yoneda\_eval}(\alpha)) =
\alpha$ uses naturality of~$\alpha$ at $f : [X \to_c A]$: evaluating
both sides at $\mathrm{id}_X$ yields
$F(f)(\alpha_X(\mathrm{id}_X)) = \alpha_A(f \circ \mathrm{id}_X)
= \alpha_A(f)$.

The mechanised round-trip laws are:

\begin{lstlisting}[
  caption={Yoneda round-trip laws (\texttt{Yoneda.v}).},
  label={lst:yoneda-roundtrip}
]
Theorem yoneda_iso_l (x : lcf_obj F X) :
    yoneda_eval (yoneda_embed x) = x.

Theorem yoneda_iso_r
    (alpha : nat_trans (repr_functor X) F) :
    yoneda_embed (yoneda_eval alpha) = alpha.
\end{lstlisting}

\paragraph{Continuity.}
To promote the bijection to a CPO isomorphism, both directions must
preserve order and $\omega$-chain suprema.  Because the
\texttt{nat\_trans} record is a plain Rocq record rather than an HB
structure---to avoid universe inconsistencies
(\cref{sec:meth-design-decisions}, DD-019)---the library states
monotonicity and continuity using standalone operations
\texttt{nt\_le}, \texttt{nt\_chain}, and \texttt{nt\_sup} rather than
the standard \texttt{monotone}/\texttt{continuous} predicates.
The four properties---\texttt{yoneda\_eval\_mono},
\texttt{yoneda\_eval\_cont}, \texttt{yoneda\_embed\_mono},
\texttt{yoneda\_embed\_cont}---together with the round-trip laws
establish that the Yoneda correspondence is a CPO isomorphism:
\[
  \mathrm{Nat}(\mathrm{Hom}(X,-),\; F)
    \;\cong_{\CPO}\;
  F(X).
\]

This result validates the enriched categorical layer: the library's
definitions of locally continuous functors, natural transformations,
and their pointwise CPO structure (\cref{sec:meth-nat-trans}) are
sufficient to state and prove the enriched Yoneda lemma with full
continuity.

% ·· Recursive domain equations ·····························
\subsection{Recursive Domain Equations}
\label{sec:res-domain-equations}

% Source: theory/DomainEquations.v (~446 lines), instances/FunLift.v (~280 lines)

\Cref{sec:bg-bilimits} defines bilimits and states the
Smyth--Plotkin existence theorem (\cref{thm:smyth-plotkin});
\cref{sec:meth-domain-equations} describes the mechanisation of the
\texttt{BilimitData} record and the \texttt{bilimit\_exists} axiom.
Here we present the concrete results: the corollaries that follow from
a \texttt{BilimitData} instance and the worked example
$D_\infty \cong [D_\infty \to_c D_\infty]_\bot$.

\paragraph{ROLL/UNROLL isomorphism.}
Given an initial object~$D_0$ and an ep-pair
$D_0 \hookrightarrow F(D_0, D_0)$, the axiom
\texttt{bilimit\_exists} yields a \texttt{BilimitData} record packaging
the bilimit~$D_\infty$, cone ep-pairs, and the isomorphism
witnesses.  The library extracts these as top-level definitions:

\begin{lstlisting}[
  caption={Bilimit corollaries (\texttt{DomainEquations.v}).},
  label={lst:bilimit-corollaries}
]
Definition D_inf  : C := bil_lim BD.
Definition ROLL   : hom (MF_obj D_inf D_inf) D_inf
  := bil_roll BD.
Definition UNROLL : hom D_inf (MF_obj D_inf D_inf)
  := bil_unroll BD.

Theorem bil_iso_roll_unroll :
    ROLL ⊚ UNROLL = id_mor D_inf.

Theorem bil_iso_unroll_roll :
    UNROLL ⊚ ROLL = id_mor (MF_obj D_inf D_inf).
\end{lstlisting}

These theorems follow directly from the record fields; the record
bundles the isomorphism witnesses together with their proofs,
so the corollaries are definitional projections.

\paragraph{Deflation chain.}
The cone ep-pairs yield a chain of deflations
$d_n = e^\infty_n \circ p^\infty_n$ in $\mathrm{hom}(D_\infty,
D_\infty)$.  Each~$d_n$ is below the identity (the deflation
condition), and the chain is increasing because each step can be
factored through the next approximant's retraction.  The mechanised
bilimit condition (\cref{eq:bilimit-identity}) is:

\begin{lstlisting}[
  caption={Deflation chain and bilimit identity
    (\texttt{DomainEquations.v}).},
  label={lst:defl-chain}
]
Definition bil_defl_chain :
    chain (hom D_inf D_inf) :=
  Build_chain
    (fun n => bil_cone_emb BD n ⊚ bil_cone_proj BD n)
    (fun m n Hmn =>
       bil_retract_chain_mono m n Hmn).

Lemma bil_sup_deflations :
    ⊔ bil_defl_chain = id_mor D_inf.
Proof.
  apply le_antisym.
  - apply sup_least; intro n.
    exact (bil_cone_deflation BD n).
  - apply bil_id_is_lub.
    intro n.
    exact (sup_upper bil_defl_chain n).
Qed.
\end{lstlisting}

The chain packages each deflation
$e^\infty_n \circ p^\infty_n$ with the monotonicity witness
built from \texttt{bil\_retract\_chain\_mono} (proved by induction
on \texttt{m $\le$ n}).  The identity proof uses antisymmetry: the
$\sqsubseteq$ direction follows because each deflation is below the
identity (\texttt{bil\_cone\_deflation}); the $\sqsupseteq$ direction
applies the record's \texttt{bil\_id\_is\_lub} field, which states
that the identity is the \emph{least} upper bound of the deflation chain.

\paragraph{Concrete instance: $D \cong [D \to_c D]_\bot$.}
\texttt{FunLift.v} registers \texttt{CPO.type} as a
\texttt{MixedLCFunctor} via the bifunctor
$F(A, B) = \langle [A \to_c B] \rangle_\bot$---the flat lift of the
continuous function space.  The object action \texttt{FL\_obj} and the
morphism action \texttt{FL\_mor} (the ``sandwich''
$h \mapsto g \circ h \circ f$ under the lift) are shown to satisfy
the six \texttt{IsMixed\-Locally\-Continuous} axioms: identity,
composition, separate monotonicity, and separate continuity.  With
this instance in place, instantiating \texttt{bilimit\_exists} with
$D_0$ the unit CPO and the canonical ep-pair yields the recursive
domain
\[
  D_\infty \cong \langle [D_\infty \to_c D_\infty] \rangle_\bot,
\]
the classical ``universal domain'' of self-applicable functions.  The
\texttt{ROLL} map wraps a continuous function into the recursive type;
\texttt{UNROLL} unwraps it, as described in \cref{sec:bg-bilimits}.

% ── Worked examples ───────────────────────────────────────────────
\section{Worked Examples}
\label{sec:res-examples}

The library ships four example files totalling approximately 1{,}000
lines of Rocq that serve as both integration tests and usage tutorials.
Each file exercises a different layer of the development, from basic
CPO constructions through PCF evaluation to recursive domain equations.

\paragraph{PCF factorial.}
The file \texttt{pcf\_examples.v} demonstrates the full verification
pipeline on a concrete recursive program.
\Cref{lst:factorial} defines a factorial function in the PCF object
language using the \texttt{FIX} combinator for general recursion,
\texttt{LET} bindings for intermediate computations, and de~Bruijn
indices (\texttt{ZVAR}, \texttt{SVAR}) for variable references.

\begin{lstlisting}[
  caption={Factorial in PCF with big-step evaluation.},
  label={lst:factorial}]
Definition factorial : CValue (Nat →ₜ Nat) :=
  FIX Nat Nat
    (LET (GT (VAR ZVAR) (NLIT 0))
        (IFB (VAR ZVAR)
            (LET (OP Nat.sub (VAR (SVAR ZVAR)) (NLIT 1))
                 (LET (APP (VAR (SVAR (SVAR (SVAR ZVAR))))
                           (VAR ZVAR))
                    (OP Nat.mul (VAR ZVAR)
                         (VAR (SVAR (SVAR (SVAR ZVAR)))))
                 ))
            (VAL (NLIT 1)))).

Example eval_factorial_5 :
    APP factorial (NLIT 5) ⇓ NLIT 120.
Proof. pcf_eval. Qed.
\end{lstlisting}

The big-step evaluation judgement
$\texttt{APP}\;\texttt{factorial}\;(\texttt{NLIT}\;5) \Downarrow
\texttt{NLIT}\;120$ is proved by a single invocation of
\texttt{pcf\_eval}, a recursive tactic combinator that tries each
big-step rule in sequence.
Because the factorial evaluation constitutes a derivation of the
big-step judgement, it can be fed directly to the soundness theorem
(\cref{sec:res-soundness}), which converts it into the denotational
equality
\[
  \llbracket\texttt{APP}\;\texttt{factorial}\;
  (\texttt{NLIT}\;5)\rrbracket
  \;=\;
  \texttt{Some}\;\llbracket\texttt{NLIT}\;120\rrbracket.
\]
Conversely, the convergence correspondence
(\cref{sec:res-correspondence}) gives the other direction: if the
denotation is defined (differs from $\bot$), the program must converge
operationally.
Together, the two results establish that
$\texttt{factorial}\;5$ converges if and only if its denotation is
non-$\bot$---exactly the bridge between operational and denotational
semantics that the library is designed to provide.

\paragraph{Kleene fixed point on $\mathbb{N}_\infty$.}
The file \texttt{basic\_cpos.v} exercises the core CPO layer
independently of PCF\@.
A simple but informative example is the Kleene fixed point of the
constant function $c_3 = \texttt{cont\_const}\;(\texttt{fin}\;3)$ on
the flat natural-number CPO $\mathbb{N}_\infty$.
The iteration chain is
\[
  \texttt{iter}\;c_3\;0 = \texttt{fin}\;0,\quad
  \texttt{iter}\;c_3\;1 = \texttt{fin}\;3,\quad
  \texttt{iter}\;c_3\;2 = \texttt{fin}\;3,\quad
  \ldots
\]
Each equation holds by \texttt{reflexivity}, confirming that the chain
stabilises after a single step.
The fixed point itself is computed by rewriting with the fixed-point
unfolding lemma:
\begin{lstlisting}[
  caption={Kleene fixed point computation on $\mathbb{N}_\infty$.},
  label={lst:kleene-demo}]
Definition c3 : [nat_inf →c nat_inf] :=
    cont_const (fin 3).

Example c3_fixp_val : fixp c3 = fin 3.
Proof.
  rewrite <- (fixp_is_fixedpoint c3).
  reflexivity.
Qed.
\end{lstlisting}
The proof applies \texttt{fixp\_is\_fixedpoint} to reduce
$\texttt{fixp}\;c_3$ to $c_3\;(\texttt{fixp}\;c_3)$; since $c_3$
ignores its argument and returns $\texttt{fin}\;3$, the goal closes by
computation.
This example connects directly to the adequacy proof
(\cref{sec:res-adequacy}), which relies on exactly this iteration-chain
machinery to extract witnesses from denotational fixed points.

\paragraph{Recursive domain equation.}
The file \texttt{recursive\_domain.v} instantiates the bilimit
machinery of \cref{sec:res-domain-equations} on the concrete
lifted-function-space bifunctor, solving
$D_\infty \cong [D_\infty \to_c D_\infty]_\bot$.
The initial object is the unit CPO, with a trivial EP-pair whose
embedding sends $\texttt{tt}$ to $\bot$ and whose projection collapses
everything to $\texttt{tt}$.
From this seed the tower
$D_0 = \texttt{unit},\; D_1 = [D_0 \to_c D_0]_\bot,\;
D_2 = [D_1 \to_c D_1]_\bot, \ldots$
is generated definitionally: each level computes by
\texttt{reflexivity}.
\Cref{lst:defl-completeness} shows the key completeness property:
the supremum of the deflation chain equals the identity,
witnessing $D_\infty$ as the colimit of the tower.

\begin{lstlisting}[
  caption={Tower levels and deflation-chain completeness.},
  label={lst:defl-completeness}]
Example tower_0 : mf_approx_obj D0 0 = D0.
Proof. reflexivity. Qed.

Example tower_1 :
    mf_approx_obj D0 1 = MF_obj D0 D0.
Proof. reflexivity. Qed.

Example sup_defls_eq_id :
    ⊔ (bil_defl_chain D0 ep0)
    = id_mor (D_inf D0 ep0).
Proof. exact (bil_sup_deflations D0 ep0). Qed.
\end{lstlisting}

The tower levels unfold by computation, confirming
that \texttt{mf\_approx\_obj} produces the expected iterated
application of the bifunctor.
The completeness proof appeals to \texttt{bil\_sup\_deflations}, which
extracts the identity from the bilimit record.
Together with the \texttt{ROLL}/\texttt{UNROLL} isomorphism shown in
\cref{sec:res-domain-equations}, this confirms that the bilimit axiom
produces a genuine solution.

\paragraph{Enriched-category constructions.}
The file \texttt{enriched\_usage.v} validates the categorical layer by
exercising representable functors, natural transformations, and
embedding--projection pairs at concrete types.
\Cref{lst:enriched-demo} shows two representative examples.
Vertical composition of natural transformations reduces to pointwise
enriched composition: the component
$(\beta \cdot \alpha)_A$ equals $\beta_A \circledcirc \alpha_A$ by
\texttt{reflexivity}, confirming that the abstract combinator
\texttt{nt\_vcomp} computes correctly.
The deflation-idempotence identity for EP-pairs similarly closes by
applying the corresponding library lemma, demonstrating that the
enriched infrastructure supports the concrete calculations
needed for recursive domain equations.

\begin{lstlisting}[
  caption={Vertical composition and EP-pair idempotence.},
  label={lst:enriched-demo},float=tp]
(* Vertical composition of natural transformations. *)
Example nt_vcomp_demo
    (alpha beta : nat_trans F0 F0) (A : CPO.type) :
    nt_component (nt_vcomp beta alpha) A =
    (nt_component beta A) ⊚ (nt_component alpha A).
Proof. exact (nt_vcomp_component beta alpha A). Qed.

(* Deflation idempotence for the identity EP-pair. *)
Example ep_idem_demo :
    (ep_emb ep ⊚ ep_proj ep) ⊚ (ep_emb ep ⊚ ep_proj ep)
    = ep_emb ep ⊚ ep_proj ep.
Proof. exact (ep_deflation_idempotent ep). Qed.
\end{lstlisting}

\medskip
\noindent
Taken together, these examples exercise all but the quantum layer of the
library---ordered sets, CPOs, continuous functions, fixed points, PCF
syntax and semantics, soundness, adequacy, enriched categories, and
recursive domains---providing confidence that the abstractions compose
correctly across the full stack.
